本文围绕专病管理中的数据割裂、风险识别滞后与院外服务不足问题，设计并实现了数据驱动的专病风险预警与智慧医患交互平台。研究工作完成了从数据接入、知识问答到风险预警与闭环干预的系统化落地，并在原型环境中验证了关键功能与工程可行性。

从方法层面看，本研究的价值在于将OCR、RAG与风险建模进行统一编排，使平台既具备智能交互能力，也具备量化预警能力；从工程层面看，版本化Prompt、可追溯日志与模块化接口设计提升了系统可维护性，为后续扩展病种与替换模型提供了基础。

后续研究将重点推进三方面工作：第一，扩大多中心脱敏数据规模，开展更严格的外部验证；第二，完善隐私保护与安全治理机制，满足更高等级合规要求；第三，探索人机协同诊疗路径，将预测结果与临床工作流深度融合，进一步提升平台在真实医疗场景中的应用价值。